\section{Условна вероятност}

\subsection{Въведение и интуиция}

Когато наблюдаваме някакъв случаен експеримент, ние работим с предварително зададени вероятности. Но какво се случва, когато получим частична информация за резултата от експеримента? Нашите очаквания, и съответно вероятностите на събитията, се променят. За да формализираме това, въвеждаме понятието за условна вероятност.

Да разгледаме един съвсем прост и класически пример – хвърляне на правилно (симетрично) зарче. Множеството от възможни изходи е $\Omega = \{1, 2, 3, 4, 5, 6\}$. При липса на каквато и да е друга информация, вероятността да се падне което и да е число $i \in \Omega$ е точно една шеста:
\[
\mathbb{P}(\{i\}) = \frac{1}{6}.
\]

Представете си сега, че имате допълнителната информация, че се е паднало нечетно число. Нека обозначим събитието „паднало се е нечетно число“ с $A = \{1, 3, 5\}$. Тогава вероятностите за отделните изходи драстично се променят. Ясно е, че ако числото е нечетно, то не може да бъде $2$. Следователно, вероятността да се е паднало $2$ при положение, че $A$ се е случило, е $0$. От друга страна, възможностите остават само три ($1, 3$ и $5$), и тъй като зарчето е симетрично, те са равновероятни. Следователно, вероятността да се е паднало $1$ при условие, че е паднало нечетно число, е вече $1/3$. 

Формално това може да се пресметне чрез отношението:
\[
\mathbb{P}(\{i\} \mid A) = \frac{\mathbb{P}(\{i\} \cap A)}{\mathbb{P}(A)}
\]
Например, ако $i = 1$, то $\{1\} \cap \{1, 3, 5\} = \{1\}$. Тогава $\mathbb{P}(\{1\} \cap A) = \frac{1}{6}$, а $\mathbb{P}(A) = \frac{3}{6} = \frac{1}{2}$. Резултатът е:
\[
\mathbb{P}(\{1\} \mid A) = \frac{1/6}{1/2} = \frac{1}{3}.
\]
Ако пък $i = 2$ (което е четно число), сечението $\{2\} \cap \{1, 3, 5\}$ е празното множество $\emptyset$, което има вероятност $0$. В този случай условната вероятност е $\frac{0}{1/2} = 0$. Този прост пример нагледно илюстрира как появата на нова информация (че се е случило събитието $A$) преизчислява нашите вероятности.

\subsection{Формална дефиниция}

\begin{defn}[условна вероятност]\label{def:cond3}
Нека е дадено вероятностно пространство $(\Omega, \mathcal{F}, \mathbb{P})$ и събитие $A \in \mathcal{F}$ такова, че $\mathbb{P}(A) > 0$. \emph{Условна вероятност на събитието $B$ при условие $A$} наричаме
\[
\mathbb{P}(B \mid A) = \frac{\mathbb{P}(A \cap B)}{\mathbb{P}(A)} .
\]
\end{defn}

Изискването $\mathbb{P}(A) > 0$ е необходимо, за да можем да извършим делението. Ако събитието $A$ има нулева вероятност (мярка $0$), в определени случаи (като при формулата за пълната вероятност) за удобство можем да приемем условната вероятност за нула, за да не усложняваме записите, макар че строго погледнато тя не е дефинирана.

Важно е да се отбележи, че функцията $B \mapsto \mathbb{P}(B \mid A)$ сама по себе си задава нова вероятностна мярка върху същото измеримо пространство $(\Omega, \mathcal{F})$. Тя притежава всички свойства на една класическа вероятност.

\subsection{Пример с Тото 6 от 49}

За да видим как частичната информация променя вероятностите в един по-реалистичен случай, да разгледаме играта Тото 6 от 49. Имате пуснат фиш с една конкретна шесторка, например $\{1, 2, 3, 4, 5, 6\}$. Нека обозначим нашата шесторка с $W_1$.
Вероятността тя да бъде изтеглена първоначално е:
\[
\mathbb{P}(W_1) = \frac{1}{\binom{49}{6}} = \frac{1}{13\,983\,816} \approx \frac{1}{N}
\]
където с $N$ сме обозначили общия брой възможни шесторки.

Представете си, че докато сте слушали тегленето по радиото, сте чули само част от числата и знаете, че в печелившата шесторка със сигурност присъстват числата $1$ и $2$. Нека това събитие да го наречем $A$. Търсим новата (условна) вероятност вашата шесторка $W_1$ да е печеливша: $\mathbb{P}(W_1 \mid A)$.

Тъй като вашата шесторка $W_1 = \{1, 2, 3, 4, 5, 6\}$ съдържа числата $1$ и $2$, то събитието $W_1$ е подмножество на събитието $A$. Тогава сечението $W_1 \cap A$ е просто $W_1$.
По дефиницията за условна вероятност:
\[
\mathbb{P}(W_1 \mid A) = \frac{\mathbb{P}(W_1 \cap A)}{\mathbb{P}(A)} = \frac{\mathbb{P}(W_1)}{\mathbb{P}(A)}
\]
Вероятността на събитието $A$ (да се изтеглят $1$ и $2$) може да се намери чрез комбинаторика. Ние вече сме фиксирали двете числа (1 и 2), така че остава да изберем останалите 4 числа от възможните 47. Броят на шесторките, съдържащи $1$ и $2$, е $\binom{47}{4}$.
Следователно $\mathbb{P}(A) = \frac{\binom{47}{4}}{\binom{49}{6}}$. Като заместим в условната вероятност:
\[
\mathbb{P}(W_1 \mid A) = \frac{1 / \binom{49}{6}}{\binom{47}{4} / \binom{49}{6}} = \frac{1}{\binom{47}{4}}
\]
Стойността $\binom{47}{4}$ е около $178\,000$. Виждаме как при наличието само на тази малка частична информация вероятността за печалба се увеличава драстично – от порядъка на $1$ към $14$ милиона на $1$ към $178$ хиляди.

\subsection{Прост модел за времето (Марковски вериги)}

Условните вероятности са в основата на стохастични процеси като Марковските вериги, които моделират системи, чието бъдещо състояние зависи само от настоящето.
Нека разгледаме един изключително опростен модел за прогноза на времето, при който дните могат да бъдат само два вида: „дъжд“ (0) и „слънце“ (1).
Можем да зададем модела чрез условни вероятности, или така наречените преходни вероятности (стохастична матрица):
\begin{align*}
\mathbb{P}(\text{утре е дъжд} \mid \text{днес е дъжд}) &= p_1 \\
\mathbb{P}(\text{утре е слънце} \mid \text{днес е дъжд}) &= 1 - p_1 \\
\mathbb{P}(\text{утре е слънце} \mid \text{днес е слънце}) &= p_2 \\
\mathbb{P}(\text{утре е дъжд} \mid \text{днес е слънце}) &= 1 - p_2
\end{align*}
Тук вероятностите за времето утре се преизчисляват при условие какво е времето днес. Ако пространството на елементарните събития за два последователни дни се зададе като декартово произведение $\{0, 1\} \times \{0, 1\}$, то $p_1$ е всъщност условната вероятност $\mathbb{P}(\text{утре е 0} \mid \text{днес е 0}) = \frac{\mathbb{P}(\{0, 0\})}{\mathbb{P}(\{0, 0\} \cup \{0, 1\})}$. Този пример илюстрира, че често на практика ние директно моделираме и задаваме условните вероятности за дадена система.

\subsection{Пример с гласоподаватели}

Друг интересен пример е свързан със социологически проучвания и избори.
Нека имаме общо гласоподаватели, от които:
\begin{itemize}
    \item $N_1$ гласуват за Партия 1.
    \item $N_2$ гласуват за Партия 2.
\end{itemize}
Нека измежду тях имаме определен брой млади гласоподаватели:
\begin{itemize}
    \item $m_1$ млади гласуват за Партия 1 ($m_1 \le N_1$).
    \item $m_2$ млади гласуват за Партия 2 ($m_2 \le N_2$).
\end{itemize}

Ако изберем случаен човек от \emph{цялата} популация, вероятността събитието $A_1$ („човекът гласува за Партия 1“) да се случи е:
\[
\mathbb{P}(A_1) = \frac{N_1}{N_1 + N_2}
\]
Ако обаче имаме информацията, че случайно избраният човек е млад (събитие $M$), каква е условната вероятност той да е гласувал за Партия 1?
Младите хора, гласуващи за Партия 1, са подмножество на всички млади хора. Следователно:
\[
\mathbb{P}(A_1 \mid M) = \frac{\mathbb{P}(A_1 \cap M)}{\mathbb{P}(M)} = \frac{\frac{m_1}{N_1 + N_2}}{\frac{m_1 + m_2}{N_1 + N_2}} = \frac{m_1}{m_1 + m_2}
\]
Забележете, че в крайния резултат $\frac{m_1}{m_1 + m_2}$ изобщо не участват $N_1$ и $N_2$. Тази пропорция може да бъде коренно различна от общата пропорция $N_1 / (N_1 + N_2)$. Този пример нагледно показва защо при статистически анализи трябва внимателно да се работи с условни вероятности, когато се ограничаваме до някакво подмножество (например определена възрастова група).

\section{Независимост на събития}

\subsection{Независимост на две събития}

Понятието за независимост е директно свързано с условната вероятност. 
\begin{defn}[независими събития]\label{def:independent}
Две събития $A$ и $B$ се наричат \emph{независими} (ще го бележим понякога с $A \perp B$), ако вероятността на тяхното сечение е равна на произведението от индивидуалните им вероятности:
\[
\mathbb{P}(A \cap B) = \mathbb{P}(A)\mathbb{P}(B) .
\]
\end{defn}

Защо това се нарича независимост? Ако приемем, че $\mathbb{P}(A) > 0$, можем да пресметнем условната вероятност на $B$ при условие $A$:
\[
\mathbb{P}(B \mid A) = \frac{\mathbb{P}(A \cap B)}{\mathbb{P}(A)} = \frac{\mathbb{P}(A)\mathbb{P}(B)}{\mathbb{P}(A)} = \mathbb{P}(B)
\]
Това означава, че информацията за настъпването на събитието $A$ не променя по никакъв начин първоначалната вероятност на събитието $B$. Събитието $A$ не носи никаква допълнителна информация за $B$. 

Ако $\mathbb{P}(A) = 0$, то сечението им автоматично има нулева вероятност, и равенството $\mathbb{P}(A \cap B) = \mathbb{P}(A)\mathbb{P}(B)$ е тривиално изпълнено ($0 = 0 \cdot \mathbb{P}(B)$).

\subsection{Независимост в съвкупност}

Често се налага да работим с повече от две събития, за които знаем (или допускаме в модела си), че са независими. Ако събитията са независими, вероятността на сеченията им се пресмята изключително лесно чрез произведения, което значително опростява изчисленията.

\begin{defn}[независимост в съвкупност]\label{def:mutual-independent}
Събитията $A_1, A_2, \dots, A_n$ се наричат \emph{независими в съвкупност}, ако за всяко подмножество от индекси $I \subseteq \{1, 2, \dots, n\}$ с поне два елемента е изпълнено:
\[
\mathbb{P}\left(\bigcap_{i \in I} A_i\right) = \prod_{i \in I} \mathbb{P}(A_i) .
\]
\end{defn}

Например, ако имаме три събития $A_1, A_2, A_3$, трябва да проверим четирите възможни подмножества от индекси с поне 2 елемента: $\{1, 2\}$, $\{1, 3\}$, $\{2, 3\}$ и $\{1, 2, 3\}$. Това означава, че трябва да са изпълнени едновременно следните условия:
\begin{align*}
\mathbb{P}(A_1 \cap A_2) &= \mathbb{P}(A_1)\mathbb{P}(A_2) \\
\mathbb{P}(A_1 \cap A_3) &= \mathbb{P}(A_1)\mathbb{P}(A_3) \\
\mathbb{P}(A_2 \cap A_3) &= \mathbb{P}(A_2)\mathbb{P}(A_3) \\
\mathbb{P}(A_1 \cap A_2 \cap A_3) &= \mathbb{P}(A_1)\mathbb{P}(A_2)\mathbb{P}(A_3)
\end{align*}

Полезна интуиция дава пресмятането на лице на правоъгълник: лицето е произведението от дължините на двете му страни $a$ и $b$. Тук страните са съответните „мерки“, а правоъгълникът е „сечението“. Тази хубава структура съществува само защото дължината и ширината на правоъгълника са ортогонални (независими) една спрямо друга.

\subsection{Приложение: Еднакви тиражи в Тотото}

Нека се върнем към един реален казус, възникнал в България през 2009 г., когато в два последователни тиража на Тото 6 от 49 се паднаха едни и същи числа. Хората тогава са се запитали: каква е вероятността за подобно нещо в историята на тотализатора?
Нека имаме 10 000 тиража. Търсим вероятността на събитието $A$: „паднали са се еднакви числа в поне два последователни тиража от общо 10 000 тиража“.
Можем да дефинираме $A_i$ като събитието „шесторката на $i$-тия тираж съвпада с шесторката на $(i+1)$-вия тираж“.
Тогава събитието $A$ е обединението на тези събития за $i$ от 1 до $9999$:
\[
A = \bigcup_{i=1}^{9999} A_i
\]
Изчисляването на вероятност на голямо обединение е трудно и би изисквало сложния принцип на включване и изключване. Много по-лесно е, когато работим с независими събития и техните сечения. Ето защо преминаваме към допълнението на $A$, използвайки формулите на Де Морган:
\[
\mathbb{P}(A) = 1 - \mathbb{P}(A^c) = 1 - \mathbb{P}\left(\left(\bigcup_{i=1}^{9999} A_i\right)^c\right) = 1 - \mathbb{P}\left(\bigcap_{i=1}^{9999} A_i^c\right)
\]
Събитията $A_i^c$ означават, че $i$-тият и $(i+1)$-вият тираж \emph{не} съвпадат. Може да се докаже, че ако събитията $A_i$ са независими в съвкупност, то и техните допълнения $A_i^c$ също са независими в съвкупност. Тогава вероятността на сечението се разпада на произведение:
\[
\mathbb{P}(A) = 1 - \prod_{i=1}^{9999} \mathbb{P}(A_i^c)
\]
Тъй като правилата на играта не се променят с времето, вероятностите $\mathbb{P}(A_i^c)$ са едни и същи за всички $1 \le i \le 9999$. Тоест $\mathbb{P}(A_i^c) = \mathbb{P}(A_1^c)$. Тогава получаваме:
\[
\mathbb{P}(A) = 1 - (\mathbb{P}(A_1^c))^{9999}
\]
Тъй като $\mathbb{P}(A_1)$ е $1 / N$ (където $N \approx 14$ милиона), то $\mathbb{P}(A_1^c) = 1 - 1/N$. Така формулата става:
\[
\mathbb{P}(A) = 1 - \left(1 - \frac{1}{N}\right)^{9999}
\]
Това приближение се пресмята изключително лесно и дори на ум може да се докаже (използвайки известни приближения от анализа като $(1-x)^n \approx 1-nx$ за малки $x$), че вероятността за това „чудо“ не е чак толкова колосално малка, колкото изглежда на пръв поглед, и истерията не винаги е математически обоснована.

\section{Пълна група от събития и Формула за пълната вероятност}

\subsection{Пълна група от събития (Хипотези)}

Когато разглеждаме някакво множество от елементарни събития $\Omega$, често е много естествено то да се раздели на групи въз основа на определена характеристика. Например, клиентите в магазин могат да се разделят на „жени“ и „мъже“. Гласоподавателите могат да се разделят на гласуващи за Партия 1, Партия 2, ..., Партия $N$ и негласуващи.

\begin{defn}[пълна група от събития]\label{def:partition}
Казваме, че събитията $H_1, H_2, \dots, H_n$ (или изброимо безкрайно много $H_1, H_2, \dots$) образуват \emph{пълна група от събития}, ако:
\begin{enumerate}
    \item Те са взаимно изключващи се (непресичащи се): $H_i \cap H_j = \emptyset$ за $i \neq j$.
    \item Те изчерпват цялото пространство: $\bigcup_{i} H_i = \Omega$.
\end{enumerate}
\end{defn}

Често тези събития $H_i$ се наричат \textbf{хипотези}, тъй като при наблюдаване на някакво следствие (данни), ние се питаме от коя точно хипотеза е породено то.

\subsection{Формула за пълната вероятност}

\begin{thm}[Формула за пълната вероятност]\label{thm:total-probability}
Нека е дадено вероятностно пространство $(\Omega, \mathcal{F}, \mathbb{P})$, нека $H_1, H_2, \dots, H_n$ е пълна група от събития и нека $A$ е някакво събитие. Тогава:
\[
\mathbb{P}(A) = \sum_{i=1}^n \mathbb{P}(A \cap H_i) = \sum_{i=1}^n \mathbb{P}(A \mid H_i)\mathbb{P}(H_i) .
\]
\end{thm}

Ако имаме изброимо безкрайно много хипотези, сумата просто продължава до безкрайност.
Ако случайно някое от събитията $H_i$ има вероятност нула, се възприема конвенцията, че съответното събираемо е нула (условната вероятност $\mathbb{P}(A \mid H_i)$ условно се приема за нула), за да не се усложнява сумата.

\begin{proof}
Тъй като $A = A \cap \Omega$, а $\Omega = \bigcup_{i=1}^n H_i$, по дистрибутивния закон имаме:
\[
A = A \cap \left( \bigcup_{i=1}^n H_i \right) = \bigcup_{i=1}^n (A \cap H_i)
\]
Тъй като $H_i$ са взаимно изключващи се, сеченията $A \cap H_i$ също са взаимно изключващи се. Следователно вероятността на тяхното обединение е сумата от вероятностите им:
\[
\mathbb{P}(A) = \sum_{i=1}^n \mathbb{P}(A \cap H_i)
\]
Като приложим дефиницията за условна вероятност $\mathbb{P}(A \cap H_i) = \mathbb{P}(A \mid H_i)\mathbb{P}(H_i)$, получаваме желаната формула.
\end{proof}

\textbf{Пример с болница:}
Представете си човек, който влиза в тежко състояние в болница. Каква е вероятността човекът да почине (събитие $A$)? Често тази обща статистика липсва или е трудна за намиране директно, но лесно може да се изведе. Нека $H_i$ са различните възможни тежки болести (хипотези). От медицинската статистика лекарите знаят честотата (вероятността) човек да се разболее от съответната болест $\mathbb{P}(H_i)$, както и смъртността при всяка болест $\mathbb{P}(A \mid H_i)$. Тогава общата вероятност се пресмята лесно чрез формулата за пълната вероятност, сумирайки по всички болести.

\section{Формула на Бейс}

Формулата на Бейс (от името на Томас Бейс) е едно от най-важните твърдения в теорията на вероятностите, стоящо в основата на Бейсовата статистика и машинното самообучение. Тя отговаря на въпроса: \emph{как да обновим вероятностите на нашите хипотези $H_k$, след като сме наблюдавали събитие (данни) $A$?}

\begin{thm}[Формула на Бейс]\label{thm:bayes}
Нека $H_1, H_2, \dots$ е пълна група от събития (хипотези) и нека $A$ е събитие с $\mathbb{P}(A) > 0$. Тогава вероятността на хипотезата $H_k$ при условие $A$ е:
\[
\mathbb{P}(H_k \mid A) = \frac{\mathbb{P}(A \mid H_k)\mathbb{P}(H_k)}{\sum_{i} \mathbb{P}(A \mid H_i)\mathbb{P}(H_i)} .
\]
\end{thm}

\begin{proof}
По дефиницията за условна вероятност:
\[
\mathbb{P}(H_k \mid A) = \frac{\mathbb{P}(H_k \cap A)}{\mathbb{P}(A)}
\]
Чрез условната вероятност числителят може да се запише като $\mathbb{P}(H_k \cap A) = \mathbb{P}(A \mid H_k)\mathbb{P}(H_k)$. Знаменателят $\mathbb{P}(A)$ се замества с разложението му от формулата за пълната вероятност. С това формулата е доказана.
\end{proof}

\subsection{Отношение между хипотези и обновяване на вероятностите}

Пресмятането на знаменателя във формулата на Бейс често пъти е компютърно много тежко (особено ако преминем към непрекъснати разпределения, където сумата става интеграл). Затова в практиката често се разглежда съотношението между две хипотези. 

Ако разгледаме две хипотези $H_k$ и $H_i$ при наличие на данни $D$, имаме:
\[
\frac{\mathbb{P}(H_k \mid D)}{\mathbb{P}(H_i \mid D)} = \frac{\frac{\mathbb{P}(D \mid H_k)\mathbb{P}(H_k)}{\mathbb{P}(D)}}{\frac{\mathbb{P}(D \mid H_i)\mathbb{P}(H_i)}{\mathbb{P}(D)}} = \frac{\mathbb{P}(D \mid H_k)}{\mathbb{P}(D \mid H_i)} \cdot \frac{\mathbb{P}(H_k)}{\mathbb{P}(H_i)}
\]
Това ни показва как апостериорното (обновеното) съотношение на хипотезите се получава чрез умножаване на априорното (първоначалното) им съотношение $\frac{\mathbb{P}(H_k)}{\mathbb{P}(H_i)}$ с фактора $\frac{\mathbb{P}(D \mid H_k)}{\mathbb{P}(D \mid H_i)}$. Това напълно елиминира трудната за пресмятане вероятност $\mathbb{P}(D)$ от знаменателя.

\subsection{Пример 1: Честна или измамна монета?}

Човек твърди, че хвърля честна монета.
\begin{itemize}
    \item $H_1$: монетата е честна (шанс за ези $1/2$).
    \item $H_2$: монетата е нечестна (за простота да допуснем, че шансът за ези е точно $1/3$).
\end{itemize}
Вие имате голямо доверие на този човек, затова вашето априорно (първоначално) субективно вярване е $\mathbb{P}(H_1) = 0{,}99$ и $\mathbb{P}(H_2) = 0{,}01$. Априорното съотношение е:
\[
\frac{\mathbb{P}(H_1)}{\mathbb{P}(H_2)} = \frac{0{,}99}{0{,}01} = 99
\]
Сега хвърляте монетата 100 пъти и събирате данни $D$: паднали са се 30 езита. Използвайки биномно разпределение, условните вероятности за тези данни са:
\begin{align*}
\mathbb{P}(D \mid H_1) &= \binom{100}{30} \left(\frac{1}{2}\right)^{30} \left(\frac{1}{2}\right)^{70} \\
\mathbb{P}(D \mid H_2) &= \binom{100}{30} \left(\frac{1}{3}\right)^{30} \left(\frac{2}{3}\right)^{70}
\end{align*}
Как се променя нашето съотношение след наблюдаването на данните $D$?
\[
\frac{\mathbb{P}(H_1 \mid D)}{\mathbb{P}(H_2 \mid D)} = \frac{\mathbb{P}(D \mid H_1)}{\mathbb{P}(D \mid H_2)} \cdot \frac{\mathbb{P}(H_1)}{\mathbb{P}(H_2)} = \frac{\mathbb{P}(D \mid H_1)}{\mathbb{P}(D \mid H_2)} \cdot 99
\]
Ако се пресметне точното отношение на биномните вероятности, крайният резултат се оказва приблизително $0{,}03$. Тоест, въпреки първоначалното ви огромно доверие (съотношение $99:1$ в полза на $H_1$), след тези експериментални данни нещата се обръщат напълно: вероятността монетата да е честна става много по-малка от вероятността да е нечестна. Тъй като $H_1$ и $H_2$ образуват пълна група, лесно може да се пресметнат самите вероятности от $\mathbb{P}(H_1 \mid D) + \mathbb{P}(H_2 \mid D) = 1$.

\subsection{Пример 2: Тестове на летището (COVID или тероризъм)}

Да си представим, че на летището е въведена система за бързо тестване (аларма) за COVID инфектирани (или пък терористи). 
Нека събитието $H_1$ означава, че човекът е инфектиран, а $H_2$ - че е здрав.
В общата популация болните са много рядко срещани, да кажем $1\%$:
\[
\mathbb{P}(H_1) = 0{,}01, \quad \mathbb{P}(H_2) = 0{,}99
\]
Тестът (алармата $B$) работи отлично и хваща почти всички инфектирани - с вероятност 99\%:
\[
\mathbb{P}(B \mid H_1) = 0{,}99
\]
Ако човекът е здрав, тестът правилно не звъни с вероятност 80\%. Следователно тестът греши (дава фалшива аларма за здрав човек) с вероятност 20\%:
\[
\mathbb{P}(B \mid H_2) = 0{,}20
\]

Вие сте охрана и алармата $B$ току-що звънна за даден човек. Каква е вероятността той наистина да е инфектиран? Търсим $\mathbb{P}(H_1 \mid B)$.
По формулата на Бейс:
\begin{align*}
\mathbb{P}(H_1 \mid B) &= \frac{\mathbb{P}(B \mid H_1)\mathbb{P}(H_1)}{\mathbb{P}(B \mid H_1)\mathbb{P}(H_1) + \mathbb{P}(B \mid H_2)\mathbb{P}(H_2)} \\
&= \frac{0{,}99 \times 0{,}01}{0{,}99 \times 0{,}01 + 0{,}20 \times 0{,}99} \\
&= \frac{0{,}0099}{0{,}0099 + 0{,}1980} \\
&= \frac{0{,}0099}{0{,}2079} \\
&\approx \frac{1}{21}
\end{align*}

Резултатът е поразителен: въпреки че тестът засича правилно 99\% от истинските болни, вероятността човекът, за когото е звъннала алармата, реално да е болен, е едва 1 към 21 (под 5\%). Причината е, че търсеното събитие в генералната популация е много рядко (0,01), а фалшиво положителните резултати ($0{,}20$) сред масовата част от хората ($0{,}99$) създават огромен брой „фалшиви аларми“, които заглушават истинските. Това илюстрира защо прилагането на подобни системи в реалността изисква изключително висока специфичност на теста, когато търсим нещо много рядко срещано.
